% =====================================================================
%  capitoli/01_introduzione.tex
% =====================================================================
\chapter{Introduction}
\label{chap:introduzione}


\section{Purpose and scope}
This document defines \emph{what} the glucose monitoring system must do, not
\emph{how} it will be built. It describes the functional architecture of the service
chosen for Project~1 of the course: a machine learning model for time series
forecasting, delivered as a web application (SaaS) on a secure cloud
infrastructure. The document is written at a high level: it presents actors,
requirements, use cases, functional components and system behavior, and it
leaves out implementation details and technology stacks, which belong to later phases.

The system is used by three kinds of users: the patients, the doctors
(diabetologists) who follow them, and the manager of the center that runs the
service. Requirements are prioritized using the MoSCoW scheme: \emph{Must}
(M), \emph{Should} (S), \emph{Could} (C) and \emph{Won't} (W).

\section{How the system works}
At every meal the patient declares the amount of carbohydrates they intend to eat. A machine learning model, trained offline on historical patient cohorts, uses the patient's recent CGM measurements and the declared carbohydrates to estimate the glucose level thirty minutes later. 

The patient immediately receives one of four possible outcomes. Two of them
are risk outcomes, for which a flag is also recorded for the doctor. No instant
notification is sent: the doctor sees the flag at the next periodic review.

The doctor consults the history periodically, for example once a week, and
identifies the trends of each patient. 


The center manager observes the flow of
requests and the performance of the cloud service.

The application has two main purposes. For patients, it helps manage carbohydrate intake at each meal instead of relying on a fixed written diet. For doctors, it provides a way to monitor patients over time, compare predicted and measured glucose levels, and decide whether to contact the patient or adjust the diet.

The system does not suggest insulin doses and does not make therapeutic
decisions: the model only forecasts glucose at $+30$ minutes, and any decision
remains with the doctor.
\section{Actors}
The system has three actors. Each of them accesses only the data and the
functions related to their  role. 

\begin{table}[H]
\centering
\caption{Actors of the system.}
\label{tab:actors}
\small
\begin{tabular}{>{\bfseries\raggedright\arraybackslash}p{3.9cm}>{\raggedright\arraybackslash}p{6.8cm}>{\raggedright\arraybackslash}p{4.6cm}}
\toprule
\rowcolor{accent}\tablehead{Actor} & \tablehead{What the actor does} & \tablehead{What the actor sees}\\
\midrule
Patient & Enters the carbohydrates of every meal and receives the outcome of the prediction & Only his or her own data\\
\rowcolor{softgrey}
Doctor (diabetologist) & Registers new patients with their diet, monitors their glucose trend and records interventions & Only the patients in his or her care\\
Center manager & Monitors the overall operation of the service & Aggregated data only\\
\bottomrule
\end{tabular}
\end{table}

\section{Data acquisition}
Four streams of data enter the system:
\begin{itemize}
  \item \textbf{Carbohydrates and meal announcements.} Entered manually by the patient at each meal, indicating the amount of carbohydrates (in grams) planned for the meal as well as any snacks scheduled across the 30-minute horizon.
  \item \textbf{Continuous glucose monitoring (CGM).} Measured by a wearable sensor that records interstitial glucose every 5 minutes without manual intervention, providing an autoregressive sliding window of recent readings ($t-12 \dots t$). In this project, the physical sensor is simulated by replaying patient time series from the OhioT1DM dataset.
  \item \textbf{Insulin records.} Basal rates and bolus doses logged by the patient's continuous subcutaneous insulin infusion (CSII) pump or smart pen device, capturing physiological insulin-on-board.
  \item \textbf{Prescribed diet.} Established by the diabetologist for each patient, specifying target ranges and maximum carbohydrate thresholds for each meal.
\end{itemize}
